Buildings obstruct camera lines of sight to hovering low-altitude unmanned aerial vehicles (UAVs). In the 2026 ShuWei Cup Problem B benchmark \cite{problem}, coarse alerts provide approximate horizontal locations and known altitudes. Mobile cameras can reach clear sight lines. The geometry of the resulting bearings determines how accurately two observers can estimate the horizontal position \cite{bishop2010}.

The field is a $100\text{ m}\times100\text{ m}$ square with five rectangular buildings and dispatch center $O=(50,50)$~m. Five observers move at 3~m/s and each occupies a disk of radius 1.5~m. We evaluate body clearance in vehicle-center configuration space \cite{lozano1983}. A gimbal mounted 2~m above each vehicle rotates in pan and tilt after the vehicle stops. A paired observation starts once both gimbals are ready, lasts 5~s, and is followed by a return to zero before the vehicle moves again.

Problem 1 maps T1 observability, maximizes inverse-range quality, and uses readiness time to choose among equal-quality stations. Problem 2(1) minimizes T2's horizontal localization error subject to a 30~s travel limit per vehicle. Problem 2(2) assigns five vehicles to six observation slots for T1--T3 and minimizes the latest paired start. Problem 3 finds the travel limit that preserves 95\% horizontal pointing coverage within $\pm5^\circ$, then derives the admissible odometry parameters for 80~m of travel. Problem 4 reuses the nominal plan to compare event-driven recourse with keeping the plan after a vehicle failure, a new alert, or a communication outage.

AI assistance in implementing the numerical checks, preparing the figures, and editing the manuscript is documented in \hyperref[sec:ai-usage]{the AI Tool Usage Report}.
